\documentclass[10pt]{article}
\usepackage[T1]{fontenc}
\usepackage[utf8]{inputenc}
\usepackage[french]{babel}
\usepackage{amsmath,amssymb,mathrsfs,mathpazo}
\usepackage{geometry}
\geometry{a4paper,margin=2cm}
\usepackage{graphicx,array,tabularx,tabularray,colortbl,multirow,multicol,booktabs,enumitem,pifont,ulem,fancybox,diagbox,cancel}
\usepackage{pstricks,pstricks-add,pst-plot,pst-tree,pst-node,pst-text,pst-func,pst-math,pst-bspline,pst-3dplot,pst-eucl,pst-solides3d}
\let\pscalculate\undefined
\usepackage{pst-bezier}
\usepackage{tikz,pgfplots,tkz-tab}
\usetikzlibrary{arrows,arrows.meta,patterns,calc,positioning,trees,decorations.markings,shapes.geometric,shapes.misc,shapes.arrows,intersections,decorations.pathmorphing,decorations.pathreplacing,backgrounds,shadows,babel}
\pgfplotsset{compat=1.18}
\usepgfplotslibrary{fillbetween}
\usepackage[np]{numprint}
\usepackage{listings,makecell,pas-tableur,twoopt}
\let\barre\undefined
\usepackage[pstricks]{tablvar}
\definecolor{darkgreen}{RGB}{0,100,0}
\providecommand{\red}{\color{red}}
\providecommand{\blue}{\color{blue}}
\providecommand{\green}{\color{green}}
\providecommand{\black}{\color{black}}
\providecommand{\white}{\color{white}}
\providecommand{\gray}{\color{gray}}
\providecommand{\cyan}{\color{cyan}}
\providecommand{\magenta}{\color{magenta}}
\providecommand{\yellow}{\color{yellow}}
\providecommand{\violet}{\color{violet}}
\providecommand{\orange}{\color{orange}}
\providecommand{\euro}{\texteuro}
\providecommand{\EUR}{\texteuro}
\providecommand{\ieme}{\textsuperscript{e}}
\providecommand{\degres}{\ensuremath{^\circ}}
\providecommand{\degree}{\ensuremath{^\circ}}
\providecommand{\vv}[1]{\overrightarrow{#1}}
\providecommand{\Coord}[1]{#1}
\providecommand{\decosix}{}
\providecommand{\starredbullet}{\ensuremath{\bullet}}
\pagestyle{empty}
\setlength{\parindent}{0pt}
\usepackage{adjustbox}\DeclareUnicodeCharacter{2236}{:}\DeclareUnicodeCharacter{2190}{\ensuremath{\leftarrow}}
\providecommand{\R}{\mathbb{R}}
\providecommand{\N}{\mathbb{N}}
\providecommand{\Z}{\mathbb{Z}}
\providecommand{\Q}{\mathbb{Q}}
\providecommand{\D}{\mathbb{D}}
\providecommand{\C}{\mathbb{C}}
\providecommand{\e}{{\mathrm{e}}}
\providecommand{\vect}[1]{\overrightarrow{#1}}
\providecommand{\vectt}[1]{\overrightarrow{\mathrm{#1}}}
\providecommand{\vvt}[1]{\overrightarrow{\mathrm{#1}}}
\providecommand{\np}[2][]{{#2~#1}}
\providecommand{\vv}[1]{{\overrightarrow{#1}}}
\renewcommand{\d}{{\mathrm{d}}}
\providecommand{\degres}{{^\circ}}
\providecommand{\degre}{{^\circ}}
\providecommand{\blue}{\color{blue}}
\providecommand{\red}{\color{red}}
\providecommand{\enskip}{\quad}
\providecommand{\llbracket}{\lbrack\!\lbrack}
\providecommand{\rrbracket}{\rbrack\!\rbrack}
\providecommand{\euro}{{\text{€}}}
\providecommand{\no}{{\text{n°}}}
\providecommand{\interval}[2]{{[#1;#2]}}
\providecommand{\intervFO}[2]{{[#1;#2[}}
\providecommand{\dotfill}{\ldots}

\providecommand{\ds}{\displaystyle}
\providecommand{\textding}[1]{#1}
\providecommand{\Oij}{\(\mathrm{O};\vec{\imath},\vec{\jmath}\)}
\providecommand{\Oijk}{\(\mathrm{O};\vec{\imath},\vec{\jmath},\vec{k}\)}
\providecommand{\Ouv}{\(\mathrm{O};\vec{u},\vec{v}\)}
\providecommand{\intervFF}[2]{\ensuremath{[#1;#2]}}
\renewcommand{\euro}{\ifmmode\text{\texteuro}\else\texteuro\fi}
\DeclareUnicodeCharacter{00B0}{\ifmmode^{\circ}\else\textdegree\fi}
\DeclareUnicodeCharacter{20AC}{\ifmmode\text{\texteuro}\else\texteuro\fi}
% Énoncés : Éducation nationale. Transcription et corrigés : APMEP, Denis Vergès et contributeurs. Découpage et classement : CoopMaths, Sébastien Lozano, Éric Elter et contributeurs. Adaptation de la présentation : Maths Explorer.
% CC BY-SA : https://coopmaths.fr/www/about/
% https://coopmaths.fr/alea/static/bac/2025/tex/bac_2025_06_sujet1_polynesie_2_cor.tex
\begin{document}
\noindent\textit{APMEP / CoopMaths. CC BY-SA. Présentation adaptée par Maths Explorer.}\par\medskip
\begin{enumerate}
\item Le sol est le plan $P_{0}$ d'équation $z=0$ donc l'avion touchera le sol lorsque $z=0$, c'est à dire $11-4t=0$ soit $t=\dfrac{11}{4}$.

Les coordonnées du point S sont donc : $\begin{cases} x = -11 + 5\times\dfrac{11}{4} = \dfrac{11}{4}  \\ y = -5 + \dfrac{11}{4} =-\dfrac{9}{4}\\ z = 0 \end{cases}$

S'il ne dévie pas de sa trajectoire,  les coordonnées du point S en lequel l'avion Bêta touchera le sol sont : $\left(\dfrac{11}{4};-\dfrac{9}{4};0\right)$.
\item
	\begin{enumerate}
		\item  L'avion Alpha transmet à la tour sa position en A$(-7; 1; 7)$ et sa trajectoire est dirigée par le vecteur $\vect{u}\begin{pmatrix}2\\-1\\-3\end{pmatrix}$.

Une représentation paramétrique de la droite $d_{\mathrm{A}}$ caractérisant la trajectoire de l'avion Alpha est donc :
$\begin{cases} x = -7 + 2s \\ y = 1 -s \\ z = 7 - 3s \end{cases}~\text{ où } s \text{ décrit } \R.$

		\item Les deux avions entreront  en collision si leur deux trajectoires se coupent, c'est à dire si les droites $d_{\mathrm{A}}$ et $d_{\mathrm{B}}$ sont sécantes.

Cherchons s'il existe un couple $(t~;~s)$ tel que le point de paramètre $t$ sur $d_{\mathrm{B}}$ soit confondu avec le point de paramètre $s$ sur $d_{\mathrm{A}}$, pour cela, résolvons :

\[
\begin{aligned}
	\begin{cases}
	-11+5t=-7+2s \\
	-5+t=1-s \\
	11-4t=7-3 s
	\end{cases}
	&\iff \begin{cases}
	-11+5 (6-s)=-7+2s \\
	t=6-s \\
	11-4(6-s)=7-3 s
	\end{cases}\\
	&\iff \begin{cases}
	-11+30-5s=-7+2s \\
	t=6-s \\
	11-24+4s=7-3 s
	\end{cases}\\
	&\iff \begin{cases}
	26=7s \\
	t=6-s \\
	7s=20
	\end{cases}\\
	&\iff \begin{cases}
	s=\dfrac{26}{7}\\
	t=6-s \\
	s=\dfrac{20}{7}
	\end{cases}
	\end{aligned}
	\]

Le système n'a pas de  solution, donc les deux avions ne peuvent pas rentrer en collision.
	\end{enumerate}
\item
	\begin{enumerate}
		\item L'avion Alpha passe par la position E$(-3~;~-1~;~1)$  s'il existe un réel $s$  tel que le point de paramètre $s$ sur $d_{\mathrm{A}}$ ait pour coordonnées $(-3~;~-1~;~1)$, pour cela, résolvons :

\[
\begin{aligned}
\begin{cases}
-3=-7+2s \\
-3=1-s \\
1=7-3 s
\end{cases}
	&\iff \begin{cases}
4=2s \\
-4=-s \\
-6=-3 s
\end{cases}\\
	&\iff \begin{cases}
s=2\\
s=2 \\
s=2
	\end{cases}
\end{aligned}
\]

Le point E$(-3~;~-1~;~1)$  correspond au point de paramètre $s = 2$ donc l'avion Alpha passe par le point~E.
		\item Le plan  $P_{\mathrm{E}}$ est perpendiculaire à la droite  $d_{\mathrm{A}}$ dirigée par le vecteur $\vect{u}\begin{pmatrix}2\\- 1\\- 3\end{pmatrix}$, donc $\vect{u}$ est un vecteur normal au plan $P_{\mathrm{E}}$.

Une équation cartésienne du plan $P_{\mathrm{E}}$ est donc de la forme : $2x - y - 3z + d = 0$ avec $d\in \R$.

De plus, le point E appartient au plan  $P_{\mathrm{E}}$ donc les coordonnées du point E vérifient l'équation du plan $P_{\mathrm{E}}$,
		d'où : d'où : $2x_E - y_E - 3z_E + d = 0$, soit $2 \times (-3) - (-1) - 3 \times 1 + d = 0$

on a donc $-6+1-3+d=0 \iff d = 8$

Conclusion : ine équation cartésienne du plan  $P_{\mathrm{E}}$ est bien : $2x - y - 3z + 8 = 0.$
		\item Déterminons les coordonnées du point F,  intersection du plan $P_{\mathrm{E}}$ et de la droite $d_{\mathrm{B}}$.

Les coordonnées des points d'intersection sont solutions du système

$\begin{cases}
x=-11+5t\\ y=-5+t \\ z=11-4t \\ 2x-y-3z+8 = 0
\end{cases}$

On cherche donc $t$ qui vérifie $2(-11 + 5t)- (-5 + t)-3(11 - 4t) + 8 = 0$ c'est-à-dire : 

$-22+10t +5-t -33 +12t+8=0 \iff 21t= 42 \iff t=\dfrac{42}{21} = 2$.

\smallskip

On en déduit
$\begin{cases}
x=-11+5t=-11+5\times 2 =-1\\
y= -5+t=-5+2=-3\\
z=11-4t=11-4\times 2 =11-8=3\\
\end{cases}$

Le point d'intersection du plan  $P_{\mathrm{E}}$ et de la droite $d_{\mathrm{B}}$ est le point  F$(-1~;~-3~;~3)$.
		\item Le repère est orthonormé on peut donc utiliser les coordonnées pour calculer les distances.

$\begin{aligned} \mathrm{EF} & =\sqrt{\left(x_\mathrm{F}-x_\mathrm{E}\right)^2+\left(y_\mathrm{F}-y_\mathrm{E}\right)^2+\left(z_\mathrm{F}-z_\mathrm{E}\right)^2}\\
\mathrm{EF} &=\sqrt{\left(-1-(-3)\right)^2+\left(-3-(-1)\right)^2+\left(3-1\right)^2}\\
	\mathrm{EF}& =\sqrt{\left(2\right)^2+\left(-2\right)^2+\left(2\right)^2}\\
	\mathrm{EF}&=\sqrt{4+4+4}\\
	\mathrm{EF}&= \sqrt{12}=2\sqrt{3}
\end{aligned}$

d'où EF$\approx \np[km]{3,4641}$ soit EF$\approx \np[m]{3646}$ à 1~m près.
	\end{enumerate}
	\item Si les avions Alpha et Bêta sont respectivement en E et F au même instant, la distance les séparant est de $\np[m]{3646}$ soit environ $\np{1215}$ milles nautiques.

Or $\np{1215}<\np{1852}$ donc leur distance de sécurité n'est pas respectée.
\end{enumerate}

\end{document}
